\subsection{Case Study 1 --- DC Motor Electromechanical System}

The direct current (DC) motor represents one of the most fundamental and widely studied electromechanical systems in control engineering. Understanding the DC motor thoroughly provides a foundation for analyzing more complex systems because it clearly demonstrates how electrical and mechanical domains interact through energy conversion. A DC motor converts electrical energy into mechanical rotational energy through the interaction of magnetic fields and current-carrying conductors. This conversion process creates a coupled system where the electrical behavior affects mechanical motion, and mechanical motion in turn influences the electrical behavior through a phenomenon called back electromotive force (back EMF). The analysis we develop in this section will illustrate the general approach to modeling coupled systems, techniques that extend directly to hydraulic, thermal, and pneumatic systems encountered throughout control systems practice.

The physical structure of a DC motor consists of two primary components: a stationary part called the stator and a rotating part called the rotor or armature. The stator produces a constant magnetic field, typically created by permanent magnets or field windings carrying direct current. The rotor carries armature windings through which current flows, and when this current interacts with the stator's magnetic field, a Lorentz force arises that produces torque on the rotor. This torque causes the rotor to rotate, and as it spins, the changing magnetic flux through the armature windings induces a voltage that opposes the applied voltage, hence the name back EMF. The complete dynamics of this system emerge from the interaction of electrical energy storage in the motor's inductance, mechanical energy storage in the rotor's rotational inertia, and energy dissipation through electrical resistance and mechanical friction.

To develop a complete mathematical model, we must capture both the electrical dynamics governing current flow and the mechanical dynamics governing rotational motion. We begin by defining the key variables and parameters that describe the system. Let $v(t)$ denote the applied armature voltage measured in volts, $i(t)$ represent the armature current measured in amperes, and $\theta(t)$ represent the angular position of the rotor measured in radians. The angular velocity $\omega(t)$ is the time derivative of position, so $\omega(t) = \dot{\theta}(t)$, and angular acceleration $\alpha(t)$ is the second derivative, $\alpha(t) = \ddot{\theta}(t)$. The torque produced by the motor is denoted $\tau_m(t)$ and the load torque applied to the motor shaft is $\tau_L(t)$. For the electrical subsystem, the armature resistance $R$ (measured in ohms) represents energy dissipation in the windings, the armature inductance $L$ (measured in henries) represents energy storage in the magnetic field, and the back EMF constant $K_e$ (measured in volts per radian per second) relates rotational velocity to induced voltage. For the mechanical subsystem, the rotor's moment of inertia $J$ (measured in kilogram meter squared) represents energy storage in rotational motion, the viscous damping coefficient $B$ (measured in newton meter second per radian) represents energy dissipation through friction, and the gear ratio $N$ captures the speed reduction and torque multiplication provided by any gearbox connected to the motor.

The coupling between electrical and mechanical domains occurs through two fundamental relationships. The first coupling relationship describes how electrical energy converts to mechanical energy: when current flows through the armature in the presence of the stator's magnetic field, a torque is produced according to the equation $\tau_m(t) = K_t i(t)$, where $K_t$ is the torque constant measured in newton meters per ampere. For a motor operating in its linear region, this relationship holds constant, and $K_t$ depends only on the physical construction of the motor, particularly the strength of the magnetic field and the geometry of the windings. The second coupling relationship describes the converse effect: when the rotor rotates, the changing magnetic flux through the armature windings induces a back EMF voltage $e_b(t)$ according to $e_b(t) = K_e \omega(t)$. This induced voltage opposes the applied voltage and therefore affects the current flow. For most DC motors, the torque constant and back EMF constant are numerically equal when expressed in consistent units, a consequence of the conservation of energy in the electromechanical conversion process. This equality means $K_t = K_e$ in SI units, and we can denote this common coupling constant simply as $K$ for the remainder of our analysis.

We now derive the complete system dynamics using the energy method, which provides an elegant approach that emphasizes the physical principles underlying the system's behavior. The total energy stored in the DC motor system consists of electrical energy stored in the armature inductance and mechanical energy stored in the rotational inertia. The electrical energy stored in the inductor is given by $W_e = \frac{1}{2} L i^2$, representing the energy in the magnetic field created by the current flow. The mechanical energy stored in the rotating mass is given by $W_m = \frac{1}{2} J \omega^2$, representing the kinetic energy of rotation. The total stored energy is the sum $W_{total} = W_e + W_m = \frac{1}{2} L i^2 + \frac{1}{2} J \omega^2$.

To derive the equations of motion, we apply the principle of power balance, which states that the rate at which energy is supplied to the system equals the rate of energy storage plus the rate of energy dissipation. The electrical power supplied to the motor is $P_{in} = v i$, the product of applied voltage and current. This power input must account for energy storage in the inductor, energy storage in the inertia, energy dissipation in electrical resistance, and energy dissipation in mechanical friction. The rate of electrical energy storage is $\frac{dW_e}{dt} = L i \frac{di}{dt}$, found by differentiating the stored electrical energy with respect to time. The rate of mechanical energy storage is $\frac{dW_m}{dt} = J \omega \frac{d\omega}{dt}$, found similarly for the mechanical case. The electrical power dissipation is $P_{R} = i^2 R$, representing joule heating in the armature windings. The mechanical power dissipation is $P_{B} = B \omega^2$, representing viscous friction losses.

Equating the input power to the sum of storage and dissipation rates gives $v i = L i \frac{di}{dt} + J \omega \frac{d\omega}{dt} + i^2 R + B \omega^2$. Dividing through by $i$ (assuming $i \neq 0$) yields $v = L \frac{di}{dt} + i R + \frac{J \omega}{i} \frac{d\omega}{dt} + \frac{B \omega^2}{i}$. This expression is not yet in standard form because of the coupling terms involving products of electrical and mechanical variables. Recognizing that the mechanical power delivered to the rotor comes from the electromechanical conversion, we identify that the conversion term contributes power $P_{conv} = \tau_m \omega = K i \omega$. The electrical equation then becomes $v = L \frac{di}{dt} + i R + K \omega$, and the mechanical equation becomes $K i = J \frac{d\omega}{dt} + B \omega + \tau_L$. These equations clearly separate into electrical and mechanical dynamics coupled through the $K\omega$ and $Ki$ terms.

The energy method derivation confirms our physical understanding, but we can also derive the same equations more directly using circuit analysis and Newton's second law for rotation. For the electrical subsystem, we apply Kirchhoff's voltage law around the armature circuit, which states that the sum of voltage drops around any closed loop equals zero. Starting from the positive terminal of the voltage source and traversing the loop, we encounter the source voltage $v(t)$, then the voltage drop across the armature resistance $iR$, the voltage drop across the armature inductance $L \frac{di}{dt}$, and finally the back EMF $e_b(t) = K\omega(t)$ which opposes the current flow. Applying Kirchhoff's voltage law gives $v(t) - i(t)R - L \frac{di(t)}{dt} - K\omega(t) = 0$, which rearranges to the electrical equation $L \frac{di(t)}{dt} + Ri(t) + K\omega(t) = v(t)$.

For the mechanical subsystem, we apply Newton's second law for rotational motion, which states that the net torque equals the moment of inertia times angular acceleration: $\sum \tau = J \alpha$. The torques acting on the rotor include the electromagnetic torque produced by the motor $\tau_m = K i$, the viscous damping torque opposing motion $\tau_B = B \omega$, and any load torque applied to the shaft $\tau_L$. Taking the direction of motor rotation as positive, the electromagnetic torque acts in the positive direction, while damping and load torques act in the negative direction. Newton's second law gives $J \frac{d\omega(t)}{dt} = K i(t) - B \omega(t) - \tau_L(t)$, which rearranges to the mechanical equation $J \frac{d\omega(t)}{dt} + B \omega(t) = K i(t) - \tau_L(t)$. When the motor operates without an external load, $\tau_L = 0$ and the equation simplifies to $J \frac{d\omega(t)}{dt} + B \omega(t) = K i(t)$.

These two first-order differential equations, one electrical and one mechanical, together describe the complete dynamics of the DC motor system. Written together, they form a coupled system:

$$
\begin{cases}
L \frac{di(t)}{dt} + Ri(t) + K\omega(t) = v(t) \\
J \frac{d\omega(t)}{dt} + B \omega(t) = K i(t) - \tau_L(t)
\end{cases}
$$

The coupling is evident: current affects angular velocity through the $K i$ term in the mechanical equation, and angular velocity affects current through the $K\omega$ term in the electrical equation. This bidirectional coupling means the system cannot be analyzed by considering electrical and mechanical parts independently; we must solve the complete system simultaneously.

\begin{center}
\begin{tabular}{|l|c|c|c|}
\hline
\textbf{Parameter} & \textbf{Symbol} & \textbf{Units} & \textbf{Physical Meaning} \\
\hline
Armature Resistance & $R$ & $\Omega$ & Energy dissipation in windings \\
\hline
Armature Inductance & $L$ & H & Energy storage in magnetic field \\
\hline
Moment of Inertia & $J$ & kg·m$^2$ & Energy storage in rotation \\
\hline
Viscous Damping & $B$ & N·m·s/rad & Energy dissipation in friction \\
\hline
Coupling Constant & $K$ & V·s/rad or N·m/A & Electromechanical conversion \\
\hline
Gear Ratio & $N$ & dimensionless & Speed reduction, torque multiplication \\
\hline
\end{tabular}
\end{center}

The gear ratio $N$ deserves special attention because it significantly affects the effective dynamics when a gearbox connects the motor to its load. A gearbox with ratio $N:1$ means that the motor shaft rotates $N$ times for each rotation of the output shaft, so the output angular velocity is $\omega_{out} = \omega_m / N$ where $\omega_m$ is the motor angular velocity. Conversely, torque is multiplied by the gear ratio, so the torque available at the output is $\tau_{out} = N \tau_m$. When analyzing a motor with a geared load, we can either model the load dynamics referred to the motor shaft (dividing load inertia by $N^2$ and load torque by $N$) or model the motor dynamics referred to the output shaft (multiplying motor inertia by $N^2$ and torques by $N$). The factor of $N^2$ for inertia arises because kinetic energy is proportional to $J\omega^2$, and since $\omega$ scales by $1/N$, the inertia must scale by $N^2$ to preserve the same energy for a given output motion.

To include gear reduction in our model, we consider a load with its own moment of inertia $J_L$ referred to the motor shaft. The total inertia seen by the motor becomes $J_{total} = J_m + J_L / N^2$, where $J_m$ is the motor's own rotor inertia. Similarly, a load torque $\tau_{L,out}$ at the output shaft appears as a reflected torque $\tau_L = \tau_{L,out} / N$ at the motor shaft. The modified mechanical equation becomes $J_{total} \frac{d\omega_m(t)}{dt} + B \omega_m(t) = K i(t) - \frac{\tau_{L,out}(t)}{N}$.

For control system design, we typically linearize the system model around a steady-state operating point. Linearization allows us to use the powerful tools of linear systems theory, including transfer functions, frequency response, and pole-zero analysis. We define small-signal perturbations about the operating point: $i(t) = I_0 + \delta i(t)$, $\omega(t) = \Omega_0 + \delta\omega(t)$, $v(t) = V_0 + \delta v(t)$, and $\tau_L(t) = \tau_{L0} + \delta\tau_L(t)$, where subscript 0 denotes the steady-state value and $\delta$ denotes the small perturbation.

At steady state, all time derivatives are zero, so the operating point satisfies $V_0 = R I_0 + K \Omega_0$ and $K I_0 = B \Omega_0 + \tau_{L0}$. These algebraic equations define the relationship between applied voltage, load torque, and steady-state speed. For a motor driving a constant load torque, increasing the applied voltage increases the steady-state speed, while increasing the load torque decreases the steady-state speed.

To linearize, we substitute the perturbed quantities into the differential equations and retain only terms linear in the perturbations. The electrical equation becomes $L \frac{d}{dt}(I_0 + \delta i) + R(I_0 + \delta i) + K(\Omega_0 + \delta\omega) = V_0 + \delta v$. Expanding and using the steady-state relationship $V_0 = R I_0 + K \Omega_0$, the constant terms cancel, leaving $L \frac{d\delta i}{dt} + R \delta i + K \delta\omega = \delta v$. Similarly, the mechanical equation becomes $J \frac{d}{dt}(\Omega_0 + \delta\omega) + B(\Omega_0 + \delta\omega) = K(I_0 + \delta i) - (\tau_{L0} + \delta\tau_L)$. Using the steady-state relationship $K I_0 = B \Omega_0 + \tau_{L0}$, we obtain $J \frac{d\delta\omega}{dt} + B \delta\omega = K \delta i - \delta\tau_L$.

The linearized equations are identical in form to the original equations but now describe small deviations from the operating point. This means the linearized model has the same structure as the original nonlinear model, but we have eliminated the constant terms that would otherwise complicate Laplace transform analysis. The linearized system represents the incremental behavior around the operating point, which is precisely what we need for controller design since controllers are designed to regulate deviations from a desired operating point.

To obtain the transfer function representation, we apply the Laplace transform to the linearized equations with zero initial conditions. The Laplace transform of the electrical equation gives $L s I(s) + R I(s) + K \Omega(s) = V(s)$, where $I(s)$, $\Omega(s)$, and $V(s)$ denote the Laplace transforms of the perturbation signals. The Laplace transform of the mechanical equation gives $J s \Omega(s) + B \Omega(s) = K I(s) - \tau_L(s)$. We can solve these algebraic equations for the transfer functions relating outputs to inputs.

For the case of no load torque perturbation ($\tau_L = 0$), we derive the transfer function from voltage input to angular velocity output. From the electrical equation, we have $I(s) = \frac{V(s) - K \Omega(s)}{Ls + R}$. Substituting this into the mechanical equation gives $J s \Omega(s) + B \Omega(s) = K \frac{V(s) - K \Omega(s)}{Ls + R}$. Multiplying both sides by $Ls + R$ yields $(J s + B)(Ls + R) \Omega(s) = K V(s) - K^2 \Omega(s)$. Bringing the $K^2 \Omega(s)$ term to the left side gives $[(J s + B)(Ls + R) + K^2] \Omega(s) = K V(s)$. Solving for $\Omega(s)/V(s)$ results in:

$$
G(s) = \frac{\Omega(s)}{V(s)} = \frac{K}{(J s + B)(L s + R) + K^2}
$$

This transfer function captures the complete dynamics of the DC motor from voltage input to velocity output. The denominator polynomial is second-order, indicating that the motor exhibits second-order system behavior with two energy storage elements (inductance and inertia). Expanding the denominator gives $G(s) = \frac{K}{LJ s^2 + (RJ + BL) s + (BR + K^2)}$. The characteristic equation $LJ s^2 + (RJ + BL)s + (BR + K^2) = 0$ determines the system's natural frequencies and damping behavior.

The transfer function from voltage to current is also useful for analysis. Setting $\tau_L = 0$ and solving for $I(s)/V(s)$, we find:

$$
\frac{I(s)}{V(s)} = \frac{J s + B}{(J s + B)(L s + R) + K^2}
$$

This transfer function shows that current responds more directly to voltage at high frequencies (where the $Js$ term dominates) and more indirectly at low frequencies (where the motor's back EMF provides feedback). The presence of $K^2$ in the denominator reflects the electromechanical coupling and shows that the back EMF acts like an additional damping term in the electrical subsystem.

For comprehensive system analysis, we also derive the state-space representation, which provides an alternative description useful for modern control design and for systems with multiple inputs and outputs. Choosing the state variables as armature current and angular velocity, $\mathbf{x} = \begin{bmatrix} i \\ \omega \end{bmatrix}$, with input $u = v$ (and load torque treated as a disturbance input), the state-space equations are:

$$
\dot{\mathbf{x}} = \begin{bmatrix}
-\frac{R}{L} & -\frac{K}{L} \\
\frac{K}{J} & -\frac{B}{J}
\end{bmatrix} \mathbf{x} + \begin{bmatrix}
\frac{1}{L} & 0 \\
0 & -\frac{1}{J}
\end{bmatrix} \begin{bmatrix} v \\ \tau_L \end{bmatrix}
$$

The output equation depends on what we wish to measure. If we measure angular velocity, $\mathbf{y} = \begin{bmatrix} 0 & 1 \end{bmatrix} \mathbf{x}$. If we measure position, we add an integrator since $\theta = \int \omega dt$, giving a third state variable. The state-space representation makes clear the bidirectional coupling: the $A$ matrix has off-diagonal terms $K/L$ and $K/J$ that couple current and velocity dynamics in both directions.

We now demonstrate parameter identification from datasheet information, an essential practical skill. Consider a representative DC motor with the following datasheet specifications: nominal voltage of 24 V, no-load speed of 3000 revolutions per minute, rated current of 2.5 A, rated torque of 0.1 N·m, armature resistance of 0.8 $\Omega$, and armature inductance of 1.5 mH. From these specifications, we can determine the coupling constant $K$ and estimate the inertia and damping parameters.

First, convert the no-load speed to radians per second: 3000 rpm equals $3000 \times 2\pi / 60 = 314.16$ rad/s. At no load, the motor produces only enough torque to overcome friction and windage losses, which we assume are small compared to rated torque. The back EMF at no-load speed is essentially equal to the nominal voltage (neglecting the small voltage drop across resistance), so $K = V_{nom} / \omega_{no-load} = 24 / 314.16 = 0.0764$ V·s/rad. Alternatively, using the rated condition, we know the back EMF at rated speed plus the resistive drop equals rated voltage: $V_{rated} = K \omega_{rated} + I_{rated} R$. Solving for $\omega_{rated}$ gives $\omega_{rated} = (V_{rated} - I_{rated} R) / K$. However, we need $\omega_{rated}$ to find $K$, creating a circular dependency. A better approach uses the torque constant directly: since $K_t = \tau_{rated} / I_{rated}$ for SI units, we find $K_t = 0.1 / 2.5 = 0.04$ N·m/A. This differs from our previous estimate because the rated speed is less than no-load speed due to the voltage drop under load. Using $K = K_t = 0.04$ V·s/rad gives consistency between electrical and mechanical specifications.

The back EMF at rated speed should satisfy $V_{rated} = K \omega_{rated} + I_{rated} R$, so $24 = 0.04 \omega_{rated} + 2.5 \times 0.8 = 0.04 \omega_{rated} + 2$. Solving gives $\omega_{rated} = (24 - 2) / 0.04 = 550$ rad/s, which equals $550 \times 60 / 2\pi = 5254$ rpm. This seems high compared to typical motors, suggesting either that the datasheet specifications are inconsistent or that the rated speed is specified at a different operating point. In practice, we would use the torque constant $K = 0.04$ as the primary parameter and estimate the rated speed from the motor's speed-torque characteristic.

To estimate the viscous damping coefficient $B$, we use the fact that at no load, the motor torque equals the friction torque. If we assume a typical friction torque of 2\% of rated torque, then $B = \tau_f / \omega_{no-load} = 0.002 / 314.16 = 6.37 \times 10^{-6}$ N·m·s/rad. The moment of inertia is not typically listed on datasheets and must be obtained from the manufacturer or measured experimentally. A typical value for a small DC motor might be $J = 1 \times 10^{-5}$ kg·m$^2$.

With identified parameters $R = 0.8 \Omega$, $L = 0.0015$ H, $K = 0.04$ V·s/rad, $J = 1 \times 10^{-5}$ kg·m$^2$, and $B = 6.37 \times 10^{-6}$ N·m·s/rad, we can compute the poles of the system. The characteristic equation is $LJ s^2 + (RJ + BL)s + (BR + K^2) = 0$. Substituting numerical values gives $(0.0015)(10^{-5}) s^2 + [(0.8)(10^{-5}) + (6.37 \times 10^{-6})(0.0015)] s + [(6.37 \times 10^{-6})(0.8) + (0.04)^2] = 0$, which simplifies to $1.5 \times 10^{-8} s^2 + 8.01 \times 10^{-6} s + 0.00161 = 0$. Dividing by $1.5 \times 10^{-8}$ gives $s^2 + 534 s + 107333 = 0$. The poles are at $s = -267 \pm j\sqrt{107333 - 267^2} = -267 \pm j\sqrt{107333 - 71289} = -267 \pm j\sqrt{36044} \approx -267 \pm j190$. The damping ratio is $\zeta = 267 / \sqrt{267^2 + 190^2} = 267 / 328 \approx 0.81$, indicating a well-damped system. The natural frequency is $\omega_n = \sqrt{267^2 + 190^2} = 328$ rad/s or about 52 Hz.

This example illustrates how datasheet specifications translate into system parameters that determine dynamic behavior. The electrical time constant $\tau_e = L/R = 0.0015/0.8 = 1.875$ ms is much smaller than the mechanical time constant $\tau_m = J/B = 10^{-5}/6.37 \times 10^{-6} = 1.57$ s. This separation of time scales (about three orders of magnitude) simplifies analysis since the electrical dynamics are much faster than mechanical dynamics. In many applications, we can approximate the electrical subsystem as instantaneous, setting $L = 0$ and replacing the electrical equation with $Ri + K\omega = v$, which reduces the system to first-order mechanical dynamics with an effective time constant determined by $J$, $B$, and $K^2/R$.

To validate the model against expected behavior, we examine several characteristic responses. Under a step voltage input with no load, the velocity should rise exponentially toward its steady-state value with a time constant determined by the system poles. The steady-state velocity is $\omega_{ss} = V/K$, which follows from setting derivatives to zero in the linearized equation: $K\omega = V$ (neglecting resistive drop for small $R$). More accurately, $\omega_{ss} = (V - I_{ss}R)/K$, and since $I_{ss} = B\omega_{ss}/K$ at no load, we have $\omega_{ss} = V / (K + B R / K)$. For our parameters, this correction is small since $BR/K = 6.37 \times 10^{-6} \times 0.8 / 0.04 = 0.000127$, much less than $K = 0.04$.

When a load torque is suddenly applied, the velocity drops and the current rises to provide additional torque. The steady-state velocity drop under load is $\Delta\omega = \tau_L / B$ for small loads where current doesn't significantly affect the dynamics. The current increase to support the load is $\Delta i = \tau_L / K$. The transient response shows the interplay between: electrical and mechanical subsystems when load torque increases, velocity decreases, which reduces back EMF, which allows more current to flow, which produces more torque to counteract the load. This negative feedback through the back EMF is what gives DC motors their self-regulating behavior.

We now extend our analysis to different load types, demonstrating how the dynamics change with the mechanical system attached to the motor. The simplest load is a constant torque load, such as a conveyor belt moving a constant mass against friction. For this load, $\tau_L$ is constant (or slowly varying compared to motor dynamics), and the analysis we have developed applies directly. The motor must produce additional torque to overcome the load, but the load itself does not store energy or introduce additional dynamics.

A more complex load is an inertial load, such as a flywheel or a vehicle accelerating from rest. This load is characterized purely by its moment of inertia, with $\tau_L = J_L \dot{\omega}_{load}$. Referred to the motor shaft with gear ratio $N$, the effective load torque is $\tau_L = (J_L / N^2) \dot{\omega}_m$. The total inertia becomes $J_{total} = J_m + J_L / N^2$, and the dynamics follow the same equations with $J$ replaced by $J_{total}$. The natural frequency decreases as $\sqrt{1/J_{total}}$, so adding inertial load slows the system's response. The settling time for velocity changes increases proportionally to $\sqrt{J_{total}}$.

A spring load represents another important category, applicable to systems like valve actuators where a mechanical spring resists motion. The load torque is proportional to position: $\tau_L = k \theta = k \int \omega dt$. This introduces an additional integrator into the system, changing the transfer function order. The mechanical equation becomes $J \dot{\omega} + B\omega + k \int \omega dt = K i$. Taking the Laplace transform gives $(Js + B + k/s) \Omega(s) = K I(s)$, which shows the spring adds a pole at the origin and a zero in the right half-plane for certain parameter combinations. Spring-loaded systems tend to oscillate at a natural frequency $\omega_s = \sqrt{k/J}$ and may require careful damping to achieve satisfactory response.

A comprehensive load combining inertia, damping, and stiffness represents the most general case: $\tau_L = J_L \dot{\omega}_{load} + B_L \omega_{load} + k_L \theta_{load}$. Referred to the motor shaft with gear reduction, these parameters transform according to $J_{ref} = J_L / N^2$, $B_{ref} = B_L / N^2$, and $k_{ref} = k_L / N^2$ for torque, while position scales as $\theta_{ref} = N \theta_{load}$. The complete system then has third-order dynamics governed by the motor inertia, motor damping, and referred load parameters. Understanding these transformations allows us to predict how the overall system behavior changes with different mechanical loads and gear ratios.

The analysis developed in this section illustrates the systematic approach to electromechanical system modeling that extends throughout control systems practice. We began with physical principles, derived coupled differential equations using both energy methods and free-body diagram approaches, linearized around an operating point, and obtained both transfer function and state-space representations. We learned to extract parameters from datasheets, validated the model against expected behavior, and extended the analysis to different load types. This methodology, with its emphasis on physical understanding combined with mathematical rigor, provides the foundation for analyzing and designing control systems across all engineering domains.
